\documentclass{tpks}
\doctitle{Стандарт на разработку требований к ПО}
\docversion{2}
\docdate{13.11.2026}

\begin{document}
\makeheader*

\section{Введение}
\subsection{Цель}
Стандарт устанавливает правила подготовки и проверки спецификации требований к ПО системы автоматизации управления жизненным циклом релиза. Проверяющий использует эти правила при приёмке новой редакции СПТ.

\subsection{Область применения}
Правила распространяются на требования высокого уровня и их изменения. Системные требования задают основание для этой работы. Архитектуру ПО и требования низкого уровня разрабатывают на следующем этапе. Процедуру разработки задаёт План разработки ПО, идентификацию версий и изменение базовых данных --- План УК ПО. Основание правил --- слайды В. П. Козырева «Требования», страницы 18, 20--27, 31--33 и 78.

\subsection{Термины, определения и сокращения}
\printterms{ПО,СПТ,ВУ,НУ,Системное требование,Производное требование,Трассируемость,Верификация,Базовая версия,ЗИ,СП,УК ПО}
ID --- постоянный идентификатор требования.

\section{Способ описания конкретного требования}
\subsection{Структура записи}
\begin{tabl}{|P{3.7cm}|L|}{Карточка требования}{\textbf{Поле} & \textbf{Правило записи}}
Идентификатор и версия & Использовать два аргумента окружения req. Вид ID и изменение версии --- по Плану УК ПО. Версия документа в шапке относится ко всей СПТ. \\ \hline
Текст & Одно наблюдаемое обязательство ПО. При наличии условия сначала назвать событие, затем ответ системы. Ветвь «иначе» добавлять, только если она определяет исход того же решения. \\ \hline
Вход и выход & Назвать доступные ПО данные, их единицы, допустимые значения и наблюдаемый результат. Для события без числовых данных указать состав и состояние до и после. \\ \hline
Тип и статус & Указать функциональное требование, характеристику качества или проектное ограничение. Указать фактический статус: проект, на проверке либо утверждено. \\ \hline
Источник и связь & Поле источника содержит родительские ID либо признак производного требования с пояснением. Формат приведён ниже. \\ \hline
Способ проверки & Указать испытание, анализ, осмотр или демонстрацию, входные условия и критерий успеха. \\ \hline
Комментарий и история & Пояснения отделить от обязательного текста. В лист изменений внести дату, основание, автора, ID и новую версию изменённого требования. \\ \hline
\end{tabl}

\subsection{Правила формулировки}
Используют «ПО должно» или «система должна». Формулировка «Если …, то ПО должно …» связывает причину с наблюдаемым выходом. Для запрета указывают результат: «ПО должно отклонить запрос». Если нужно описать отсутствие воздействия, используют форму «должно не …». Слова «удобно», «быстро» и «при необходимости» заменяют проверяемым условием. Технологию указывают только как обоснованное проектное ограничение. Если значение неизвестно, записывают причину, владельца вопроса и условие снятия. До разрешения вопроса затронутую карточку сохраняют в проекте.

\clearpage
\subsection{Образец записи}
Условие примера взято из \reqref{SYS_01_05}: для регистрации версии необходимы все обязательные сведения о сборке. Автор заполняет окружение \texttt{req}; первый аргумент задаёт ID, второй --- номер редакции. Пример предназначен для иллюстрации формата, а не для включения в реестр утверждённых требований.

\noindent\textbf{Запись в исходном документе}
\begin{verbatim}
\begin{req}{SRD_01_01}{1}
Если в сообщении о сборке зарегистрированного компонента
отсутствует номер версии, то ПО должно отклонить регистрацию
версии.
\reqsource{parent=[SYS_01_05]}
\reqverify{Испытание. Передать сообщение без номера версии.
Операция отклонена, новой записи версии нет.}
\reqnote{Вход: сообщение без номера версии. Выход: отказ.
Тип: функциональное. Статус: проект.}
\end{req}
\end{verbatim}

\noindent\textbf{Та же карточка после сборки документа}
\begin{req}{SRD_01_01}{1}
Если в сообщении о сборке зарегистрированного компонента отсутствует номер версии, то ПО должно отклонить регистрацию версии.
\reqsource{parent=[SYS_01_05]}
\reqverify{Испытание. Передать сообщение без номера версии. Операция отклонена, новой записи версии нет.}
\reqnote{Вход: сообщение без номера версии. Выход: отказ. Тип: функциональное. Статус: проект.}
\end{req}

\subsection{Запись ссылок}
Источник записывают один раз командой \verb|\reqsource|. Образец поля:
\begin{quote}
\verb|need=[П-1,П-9]|
\end{quote}
Имена пишут строчными латинскими буквами. Поля разделяют точкой с запятой, значения --- запятой. Между полями допустимы пробелы и переносы строк. Значения перечисляют полностью, без диапазонов, повторов и пустых элементов. Порядок элементов списка не влияет на смысл. Каждое имя встречается в записи один раз.

Для SYS обязательны поля \texttt{need}, \texttt{use} и \texttt{draft} в этом порядке. Первые два поля указывают, какую потребность удовлетворяет требование и в каком сценарии оно действует. Значения П-N и ВИ-N берут из таблицы потребностей и названий сценариев документа 01, где N --- положительное целое число. У требования \reqref{SYS_01_04} источник записан так:

\begin{quote}
\texttt{need=[П-1,П-9]; use=[ВИ-1]; draft=[SYS-004]}
\end{quote}

Поле draft указывает происхождение в части Б рабочих материалов команды: SYS-NNN --- трёхзначный номер чернового требования, R-NN --- двузначный номер резервной записи, Б.N --- номер раздела. Значение TEAM обозначает решение команды, основание которого фиксируют в истории документа. Так, SYS-004 указывает на черновик. В принятой спецификации регистрация имеет ID SYS_01_04.

Для SRD используют одно поле \texttt{parent}. Связь направлена от требования к ПО к каждому системному требованию, которое оно уточняет. Пример записи:
\begin{quote}
\texttt{parent=[SYS_01_04,SYS_01_05]}
\end{quote} Такая запись уместна, только если текст уточняет оба основания. Для производного требования единственная запись источника --- \texttt{derived=[yes]}. В комментарии автор указывает исходное проектное допущение, причину появления требования и кому передано обоснование для рассмотрения на системном уровне.

\begin{tabl}{|P{3.6cm}|L|L|}{Проверка трассировки}{\textbf{Проверка} & \textbf{Действие} & \textbf{Критерий}}
Формат записи & Сверить имена полей, разделители и типы значений с правилами выше. & Каждый ID указан полностью. В списках нет повторов и пустых элементов. \\ \hline
От требования к источнику & Для каждой цели найти её определение в зафиксированной редакции входного документа. & Каждой ссылке соответствует ровно одна цель нужного типа. Версию документа можно восстановить по записи УК. \\ \hline
От источника к требованиям & Выбрать все карточки, в которых указан данный источник. Для сценария сверить список с указателем документа 01, для системного требования --- с объёмом текущей СПТ. & Все карточки, уточняющие источник, найдены. По каждому пробелу в покрытии принято решение. \\ \hline
Производное требование & Проверить поле derived и прочитать обоснование в комментарии. & Указаны причина, исходное допущение и передача на системный уровень. Поле parent отсутствует. \\ \hline
\end{tabl}

\section{Описание процесса}
\subsection{Окружение процесса}
\begin{center}
\setlength{\unitlength}{1mm}
\begin{picture}(155,60)
\put(41,52){\makebox(26,7){\small План разработки ПО}}
\put(94,52){\makebox(50,7){\small\shortstack{Стандарт на разработку\\требований к ПО}}}
\put(54,52){\line(0,-1){7}}\put(54,45){\line(1,0){16}}\put(70,45){\vector(0,-1){9}}
\put(119,52){\line(0,-1){7}}\put(119,45){\line(-1,0){28}}\put(91,45){\vector(0,-1){9}}
\put(56,18){\framebox(48,18){\shortstack{\textbf{Разработка}\\\textbf{требований к ПО}}}}
\put(0,37){\makebox(50,6)[l]{\small Системные требования}}
\put(0,37){\line(1,0){48}}\put(48,37){\line(0,-1){5}}\put(48,32){\vector(1,0){8}}
\put(0,27){\makebox(49,6)[l]{\small Архитектура системы}}
\put(0,27){\vector(1,0){56}}
\put(0,12){\makebox(48,10)[l]{\small\shortstack[l]{Программно-аппаратные\\интерфейсы}}}
\put(0,12){\line(1,0){48}}\put(48,12){\line(0,1){10}}\put(48,22){\vector(1,0){8}}
\put(104,31){\vector(1,0){49}}\put(107,32){\makebox(46,11)[l]{\small\shortstack[l]{Спецификация\\требований к ПО}}}
\put(104,22){\vector(1,0){49}}\put(107,11){\makebox(46,10)[l]{\small\shortstack[l]{Сообщения\\о проблемах}}}
\end{picture}
\captionof{figure}{Окружение процесса по слайду 18}
\end{center}
При повторном входе утверждённый ЗИ определяет границы изменения базовой версии. Его рассматривают отдельно от трёх исходных данных схемы.

\subsection{Критерий начала}
Автор назначен, базовые версии системных требований, архитектуры и интерфейсов определены. Автору доступны применимые редакции плана и стандарта. Для изменения существующей базовой версии назначен ЗИ с указанием изменяемой ЕК.

\subsection{Входные данные}
Системные требования задают поведение продукта. Архитектура системы распределяет функции между системой, пользователями и внешними средствами. Программно-аппаратные интерфейсы задают границы и направление обмена. В этом проекте два последних входа представлены разделом 3 документа 01. Автор записывает редакцию используемого документа перед началом работы.

\subsection{Выходные данные}
Спецификация требований к ПО содержит проверяемые требования и их связи. Сообщения о проблемах описывают найденные несоответствия входов или результатов. Например, если требуемое состояние релиза невозможно определить по данным внешней системы, автор указывает в СП исходное требование и недостающее наблюдение.

\subsection{Структура процесса}
\begin{enumerate}
\item Определить входы и выходы ПО. Для каждого события записать отправителя, данные, условия допустимости и наблюдаемый ответ. Для ошибки или задержки установить, какой результат остаётся известен.
\item Сформулировать требования к ПО как к «чёрному ящику». Описать штатные и предусмотренные нештатные режимы через результаты на границе ПО. Для каждого утверждения указать способ отличить успех от ошибки.
\item Трассировать требования к породившим их системным требованиям. Для производного требования обосновать происхождение. Построить обратную выборку и разобрать каждый случай, когда системное обязательство в объёме поставки осталось без уточнения.
\end{enumerate}

\subsection{Критерий окончания}
Независимый проверяющий сверил СПТ со входами и настоящим стандартом. Все замечания, препятствующие утверждению, закрыты. Утверждённая редакция передана в УК как базовая версия вместе с результатом проверки. Проектировщик может получить ту же редакцию по записи УК.

\subsection{Роли и квалификация участников}
\par\noindent\begin{minipage}{\linewidth}
\captionof{table}{Участники процесса разработки требований}
\begin{tabularx}{\linewidth}{|P{3.3cm}|L|L|}
\hline \textbf{Роль} & \textbf{Что участник должен уметь} & \textbf{За что отвечает} \\ \hline
Аналитик & Прочитать системное требование, таблицу обмена и сценарий. Выделить событие, условия и результат. & Готовит СПТ, устанавливает связи, обосновывает производные требования, регистрирует СП. \\ \hline
Проверяющий & Проследить связь в обоих направлениях и применить критерий проверки к входным данным. & Независимо рассматривает редакцию и записывает результат проверки и замечания. \\ \hline
Ответственный за УК & Идентифицировать и восстановить редакцию документа по записи УК. & Регистрирует утверждённую версию, сохраняет результат проверки и передаёт версию проектировщику. \\ \hline
\end{tabularx}
\end{minipage}\par
Проверку выполняет другой участник. Персональные назначения фиксируют по Плану разработки ПО.

\section{Правила составления спецификации}
\subsection{Содержание спецификации}
Автор описывает четыре группы сведений, указанные на слайде 20:
\begin{enumerate}
\item Входные и выходные данные ПО: источники, получатели, значения, единицы и условия допустимости.
\item Требования для всех режимов работы: штатное продвижение, отказ, неизвестный результат задания и восстановление после перезапуска в пределах принятых системных требований.
\item Протоколы взаимодействия ПО с аппаратурой и другими системами: последовательность сообщения и ответа, условия ошибки, повтор и ожидание. Особенности аппаратуры берут из системного окружения.
\item Ограничения на проектирование с источником и обоснованием. Например, согласованный способ аутентификации берут из архитектурного раздела системных требований.
\end{enumerate}

\subsection{Свойства отдельного требования}

\subsubsection{Корректность}
Ожидаемый результат следует из исходных данных. У производного требования проверяют обоснование.

\textit{Проверка.} Сравнить перечень обязательных полей с системным требованием о регистрации сборки.
\subsubsection{Однозначность}
Два читателя одинаково определяют событие, исполнителя и результат.

\textit{Проверка.} Для фразы «после выпуска» указать точное событие: завершение установки или подтверждение выпуска.
\subsubsection{Неделимость (атомарность)}
Карточка задаёт одно обязательство. Если одно действие можно принять отдельно от другого, их разделяют.

\textit{Проверка.} Проверить, можно ли независимо принять регистрацию версии и соблюдение предельного времени регистрации.
\subsubsection{Верифицируемость (тестируемость)}
Для заданных условий можно определить, выполнено ли требование. Метод проверки и критерий указаны в карточке.

\textit{Проверка.} Подготовить вход без номера версии. Успех проверки означает, что операция отклонена и новая версия не появилась.
\subsubsection{Трассируемость}
По источнику можно найти основание, по обратной связи --- уточняющие требования. Производное требование явно отмечено.

\textit{Проверка.} От примера SRD_01_01 перейти к SYS_01_05 и найти ту же карточку обратной выборкой.

Величины по возможности выражают числом с единицей измерения. Для входов задают допустимые диапазоны, для вычисляемых выходов --- допустимую погрешность. Проверяют совместимость с целевой аппаратурой. Детали реализации допускаются только для установленного проектного ограничения.

\subsection{Свойства спецификации в целом}

\subsubsection{Полнота}
Для каждого системного обязательства в границах СПТ определены поведение, входы, выходы и соответствующие режимы. Проверяющий разбирает пропуски обратной трассировкой.

\textit{Проверка.} Для промышленного выпуска проследить оба исхода: подтверждение до окончания окна и откат при отсутствии подтверждения.
\subsubsection{Непротиворечивость (согласованность)}
Результаты совместимы при одинаковых исходных условиях. Например, повтор установки согласуется с правилом блокировки среды.

\textit{Проверка.} Сопоставить условия повтора установки с условиями блокировки среды при неизвестном результате задания.
\subsubsection{Неизбыточность}
Каждое обязательство сформулировано один раз. В других разделах помещают ссылку на его ID.

\textit{Проверка.} Если правило авторизации требуется в двух сценариях, проверить, что оба ссылаются на его единственную карточку.
\subsubsection{Модифицируемость}
ID сохраняется. По его связям можно определить затронутые карточки, изменить каждую в одном месте и воспроизвести предыдущую редакцию.

\textit{Проверка.} Для изменившегося системного требования выбрать затронутые требования к ПО и восстановить их предыдущие редакции.


\begin{changelog}
\change{1}{13.11.2026}{Первая версия документа}{Пономаренко А. А.}{Задание к паре 31.10.2026}
\change{2}{13.11.2026}{Введены образец записи и его печатный вид, схема окружения, таблицы участников и проверок, проверки свойств СПТ}{Пономаренко А. А.}{Переработка по образцам прошлых проектов}
\end{changelog}
\end{document}
